\documentclass[11pt,a4paper]{article}

% ============================================================
% Standard English mathematical typesetting preamble
% Cayley, Collected Mathematical Papers, Vol.\ VII
% Book pages 26--50 (papers 423--433).  Scan-verified wrapper.
% ============================================================
\usepackage[utf8]{inputenc}
\usepackage[T1]{fontenc}
\usepackage[english]{babel}
\usepackage{amsmath,amssymb,amsthm}
\usepackage{geometry}
\geometry{margin=1in}
\usepackage{graphicx}
\usepackage{booktabs}
\usepackage{array}
\usepackage{fancyhdr}
\usepackage{titlesec}
\usepackage{enumitem}
\usepackage{mathtools}
\usepackage{bm}

\theoremstyle{plain}
\newtheorem{theorem}{Theorem}
\newtheorem{lemma}{Lemma}
\theoremstyle{definition}
\newtheorem{remark}{Remark}

\pagestyle{fancy}
\fancyhf{}
\fancyhead[C]{\small Cayley --- Collected Mathematical Papers, Vol.\ VII (pp.\ 26--50)}
\fancyfoot[C]{\thepage}
\renewcommand{\headrulewidth}{0.4pt}

\titleformat{\section}{\large\scshape\centering}{}{0em}{}
\titleformat{\subsection}{\normalsize\bfseries}{}{0em}{}

\newcommand{\dd}{\,d}
\DeclareMathOperator{\Jac}{Jac}

\begin{document}

% ============================================================
% PAPER 423 (continues from earlier portion of the volume)
% Book pp.\ 26--30
% ============================================================
\section{[423]}
\subsection{On the Plane Representation of a Solid Figure.}

\begin{center}
\textit{[From the Philosophical Magazine, vol.~\textsc{xli}. (1871), pp.~286--290.]}
\end{center}

\textsc{We} represent \textit{in plano} the position of a point~$P$ whose coordinates in space are
$(x,\,y,\,z)$ by drawing these coordinates, on the same scale or on different scales, and
in given directions from a fixed origin in the plane; $OM = x$, $MP' = y$, $P'P'' = z$.  But
observe that the point~$P''$ \textit{alone} does not completely represent the point~$P$; in fact
$P''$ represents a whole series of points lying in a line; any one such point is the
point whose coordinates are $Om$, $mp'$, $p'P''$.  For the complete representation of~$P$ we
require the \textit{two points} $P'$, $P''$: these might be distinguished as the projection~$P''$, and
the foot-point~$P'$.  The two points $P'$, $P''$ are obviously such that the line joining them
is in a given direction.

The preceding is, of course, the ordinary method of orthogonal projection, or
geometrical delineation of a solid figure: it may be used under various forms; for
example, the coordinates $x$, $y$, $z$ may be taken on the same scale and in directions
inclined to each other at angles of $120^\circ$ (isometrical projection); or the coordinates $x$, $y$
may be drawn on the same scale and at their actual inclination, $90^\circ$, to each other;
and the coordinate~$z$ on the same or an altered scale in any given direction; the
points $P'$ then give a true ground-plan of the solid figure, and the lengths of the
lines $P'P''$ give the altitudes of the several points~$P$: this is also a method in
ordinary use.

But it is to be observed that the points $P'$, $P''$ are both of them \textit{projections},
and that the general theory is as follows: we represent the position of the point~$P$
by means of its projections $P'$, $P''$, from two fixed points $\Omega'$, $\Omega''$ respectively; the
line joining these points passes, it is clear, through a fixed point~$\Omega$ which is the
intersection of the plane of projection by the line which joins the two points $\Omega'$, $\Omega''$.

Hence we say that a point~$P$ in space is represented \textit{in plano} by any two points
$P'$, $P''$ which are such that the line joining them passes through a fixed point~$\Omega$.
And we have thus a \textit{system of constructive geometry} which is the more simple on
account of the generality of its basis, and which is at once applicable to any of the
special projections above referred to.  I establish the fundamental notions of such a
geometry, and by way of illustration apply it to the solution of the well-known problem
of finding the lines which meet four given lines in space.

A point~$P$ (as already mentioned) is given by its projections $P'$, $P''$, which are
points such that the line joining them passes through the fixed point~$\Omega$.

A line~$L$ is given by its projections $L'$, $L''$, which are any two lines in the plane.
We speak of the point $(P',\,P'')$, meaning the point~$P$ whose projections are $P'$, $P''$;
and similarly of the line $(L',\,L'')$, meaning the line whose projections are $L'$, $L''$.

If $P'$, $P''$ coincide, then the point~$P$ is in the plane of projection; and so if
$L'$, $L''$ coincide, then the line~$L$ is in the plane of projection.

If through~$\Omega$ we draw a line meeting $L'$, $L''$ in the points $P'$, $P''$ respectively,
these are the projections of a point~$P$ on the line~$L$.  In particular the intersection
of $L'$, $L''$ (considered as two coincident points) represents the intersection of the line
$L$ with the plane of projection.

\clearpage

The line through the points $(P',\,P'')$ and $(Q',\,Q'')$ has for its projections the lines
$P'Q'$ and $P''Q''$.

Two lines $(L',\,L'')$ and $(M',\,M'')$ intersect each other if only the intersections $L'M'$
and $L''M''$ are the projections of a point~$P$---that is, if the line through the points
$L'M'$ and $L''M''$ passes through~$\Omega$.  And then clearly $P$ is the intersection of the two
lines.

A plane~$\Pi$ is conveniently given by means of its trace~$\Theta$ on the plane of
projection, and of the projections $(P',\,P'')$ of a point on the plane; or, say, by means
of the trace~$\Theta$, and of a point~$P$ on the plane.

Suppose, however, that a plane is given by means of a line~$L$ and a point~$P$
on the plane.  The trace~$\Theta$ passes through the point of intersection of the line~$L$ with
the plane of projection---that is, through the point of intersection of the projections
$L'$, $L''$.  To find another point on the trace, we have only to imagine on the line~$L$
a point~$Q$, and, joining this with~$P$, to suppose the line~$PQ$ produced to meet the
plane of projection.  The construction is obvious; but by way of illustration I give it
in full.  Through $\Omega$ draw a line meeting $L'$, $L''$ in $Q'$, $Q''$ respectively (then these are
the projections of a point~$Q$ on the line~$L$); the lines $P'Q'$ and $P''Q''$ are the
projections of the line~$PQ$, and the intersection of $P'Q'$ and $P''Q''$ is therefore the required
point on the trace~$\Theta$.

The line of intersection of two planes passes through the point of intersection of
their traces $\Theta_1$, $\Theta_2$; whence, if the planes have in common a point~$P$, the line of
intersection is the line joining $P$ with the intersection of the traces $\Theta_1$, $\Theta_2$.

In what precedes we have the solution of the following problem:---``Given a point
$P$, and two lines $L_1$, $L_2$, to find a line through~$P$ meeting the two lines $L_1$, $L_2$.''  The
required line is in fact the line of intersection of the planes $(P,\,L_1)$ and $(P,\,L_2)$; we
have seen how to construct the traces $\Theta_1$ and $\Theta_2$ of these planes respectively; and
the required line is the line joining $P$ with the intersection of $\Theta_1$ and $\Theta_2$.

I proceed now to the problem to find the two lines, each of them meeting four
given lines, $L_1$, $L_2$, $L_3$, $L_4$ (these being, of course, given by means of their projections
$(L_1',\,L_1'')$ \&c.).  The question is in effect to find on the line~$L_1$ a point~$P$ such that,
drawing from it a line to meet $L_2$, $L_3$, and also a line to meet $L_3$, $L_4$, these shall
be one and the same line.

Now, considering in the first instance $P$ as an arbitrary point on the line~$L_1$,
the line from~$P$ to meet $L_2$, $L_3$ is any line whatever meeting the lines $L_1$, $L_2$, $L_3$:
say it is a generating line of the hyperboloid whose directrices are $L_1$, $L_2$, $L_3$, or of
the hyperboloid $L_1L_2L_3$.  Hence projecting from any point $\Omega'$ whatever, the generating
lines and directrices are projected into tangents of one and the same conic.  We know
the projections $L_1'$, $L_2'$, $L_3'$ of the directrices; to find two other tangents of the conic,
we take two arbitrary positions of $P$ on the line~$L_1$, and construct as above the
projections $M'$, $N'$ of the lines from these to meet the lines $L_2$, $L_3$.  The conic is then

\clearpage

given as the conic touching the five lines $L_1'$, $L_2'$, $L_3'$, $M'$, $N'$: say this is the conic
$\Sigma'$.  Similarly, instead of $\Omega'$, considering the point $\Omega''$, we have the lines $L_1''$, $L_2''$, $L_3''$,
and the lines $M''$, $N''$, which are the other projections of the lines through the two
positions of $P$; and touching these five lines we have a conic $\Sigma''$.  Each tangent $T'$
of $\Sigma'$, combined with the corresponding tangent $T''$ of $\Sigma''$, represents a line $T$ meeting
$L_1$, $L_2$, $L_3$; to establish the correspondence, observe that, inasmuch as the line $T$ meets
$L_1$, the intersections of $T'$, $L_1'$ and of $T''$, $L_1''$ must lie in a line with $\Omega$; if $T'$ be
given, the point $(T'',\,L_1'')$ is thus uniquely determined, and therefore also $T''$ (since $L_1''$
is a tangent of $\Sigma''$); and similarly if $T''$ be given, $T'$ is uniquely determined; the
correspondence $T'$, $T''$ is thus, as it should be, a $(1,\,1)$ correspondence.

Considering in like manner the lines which meet $L_1$, $L_3$, $L_4$, we have touching
$L_1'$, $L_3'$, $L_4'$, $M'$, $N'$ a conic $\Sigma_1'$; and touching $L_1''$, $L_3''$, $L_4''$, $M''$, $N''$ a conic $\Sigma_1''$; each
tangent $T'$ of $\Sigma_1'$, combined with the corresponding tangent $T''$ of $\Sigma_1''$, represents a line
meeting $L_1$, $L_3$, $L_4$, the correspondence being a $(1,\,1)$ correspondence such as in the
former case.

The conics $\Sigma'$, $\Sigma_1'$ both touch $L_1'$, $L_3'$; hence they have in common two tangents.
Say one of these is $T' = T_1'$; the corresponding tangents $T''$ and $T_1''$ will coincide
with each other and be a common tangent of $\Sigma''$, $\Sigma_1''$ (these conics both touch $L_1''$, $L_3''$,
and have thus in common two tangents).  We have thus $T' = T_1'$, and $T'' = T_1''$, as the
projections of a line meeting $L_1$, $L_2$, $L_3$, $L_4$; and taking the other common tangents
of $\Sigma'$, $\Sigma_1'$ and of $\Sigma''$, $\Sigma_1''$, we have the projections of the other line meeting $L_1$, $L_2$, $L_3$, $L_4$.

The whole process is:---Construct $M'$, $M''$ and $N'$, $N''$ each of them the projections
of a line through a point $P$ of $L_1$, which meets $L_2$, $L_3$; and $M_1'$, $M_1''$ and $N_1'$, $N_1''$ each
of them the projections of a line through a point $P$ of $L_1$, which meets $L_3$, $L_4$; we
have then the conics
\begin{align*}
\Sigma',\ \Sigma''\quad &\text{touching}\quad L_1',\,L_2',\,L_3',\,M',\,N'\ \text{and}\ L_1'',\,L_2'',\,L_3'',\,M'',\,N''\ \text{respectively}, \\
\Sigma_1',\ \Sigma_1''\quad &\text{touching}\quad L_1',\,L_3',\,L_4',\,M_1',\,N_1'\ \text{and}\ L_1'',\,L_3'',\,L_4'',\,M_1'',\,N_1''\ \text{respectively};
\end{align*}
and then the projections of each of the required lines are $T' = T_1'$, a common tangent
of $\Sigma'$, $\Sigma_1'$, and $T'' = T_1''$, the corresponding common tangent of $\Sigma''$, $\Sigma_1''$.

It is material to remark how the construction is simplified when there is given
one of the lines, say~$M$, which meets $L_1$, $L_2$, $L_3$, $L_4$.  Here if $M$ is a common directrix
of the two hyperboloids, we may for the hyperboloids $\Sigma'$ and $\Sigma''$ consider, instead of
$L_1$, $L_2$, $L_3$ and two new generating lines, the lines $L_1$, $L_2$, $L_3$, $M$, and a single new
generating line~$N$; and similarly for the hyperboloids $\Sigma_1'$, $\Sigma_1''$ the lines $L_1$, $L_3$, $L_4$, $M$, and
a single new generating line~$N_1$.  $\Sigma'$, $\Sigma_1'$ have thus in common the three tangents
$L_1'$, $L_3'$, $M'$, and therefore only a single other common tangent, $T' = T_1'$; and similarly
$\Sigma''$, $\Sigma_1''$ have in common the three tangents $L_1''$, $L_3''$, $M''$, and therefore only a single
other common tangent, $T'' = T_1''$; and we have thus the other line cutting the four
given lines.

\clearpage

I take the opportunity of mentioning the following theorem:

``If in a given triangle we inscribe a variable triangle of given form, the envelope
of each side of the variable triangle is a conic touching the two sides (of the given
triangle) which contain the extremities of the variable side in question.''

We have thence a solution of the problem (\textit{Principia}, Book~I. Sect.~V.\ Lemma
XXVII.), in a given quadrilateral to inscribe a quadrangle of given form.  The question
in effect is: in the triangle $ABC$ to inscribe a triangle $\alpha\beta\gamma$ of given form; and in
the triangle $ADE$ a triangle $\alpha'\beta'\gamma'$ of given form, in such wise that the sides $\alpha\gamma$, $\alpha'\gamma'$
may be coincident.  The envelope of $\alpha\gamma$ is a conic touching $AD$, $AE$, and the envelope
of $\alpha'\gamma'$ a conic also touching $AD$, $AE$: there are thus two other common tangents, either
of which may be taken for the position of the side $\alpha\gamma = \alpha'\gamma'$; and the problem admits
accordingly of two solutions.

\clearpage

% ============================================================
% PAPER 424
% Book pp.\ 31--33
% ============================================================
\section{[424]}
\subsection{On the Attraction of a Terminated Straight Line.}

\begin{center}
\textit{[From the Philosophical Magazine, vol.~\textsc{xli}. (1871), pp.~358--360.]}
\end{center}

\textsc{Write} for shortness $(a,\,b,\,c\,;\,\varepsilon)$ to denote the shell included between the ellipsoids
\[
\frac{x^2}{a^2} + \frac{y^2}{b^2} + \frac{z^2}{c^2} = 1
\quad\text{and}\quad
\frac{x^2}{a^2} + \frac{y^2}{b^2} + \frac{z^2}{c^2} = (1+\varepsilon)^2,
\]
(where $\varepsilon$ is indefinitely small); then, if the ellipsoids
\[
\frac{x^2}{a^2} + \frac{y^2}{b^2} + \frac{z^2}{c^2} = 1
\quad\text{and}\quad
\frac{x^2}{a'^2} + \frac{y^2}{b'^2} + \frac{z^2}{c'^2} = 1
\]
are confocal, the attractions of the shells $(a,\,b,\,c\,;\,\varepsilon)$ and $(a',\,b',\,c'\,;\,\varepsilon)$ upon any exterior
point~$P$ are proportional to their masses.  Hence, considering a prolate spheroid of
revolution, $c=b$, the attractions of the shell $(a,\,b,\,b\,;\,\varepsilon)$ will be proportional to those
of the shell $\bigl(\sqrt{a^2-h},\ \sqrt{b^2-h},\ \sqrt{b^2-h}\,;\,\varepsilon\bigr)$; or if, as usual, $b^2 = a^2(1-e^2)$, then, if $h$ increases
and becomes ultimately equal to $b^2$, to those of the shell $(ae,\,0,\,0\,;\,\varepsilon)$; viz.\ this last
is the portion of the axis of $x$ included between the limits $x = -ae$, $x = +ae$; or say
it is the terminated line $x = \pm\, ae$; and I say that the mass is distributed over this line
uniformly.

To see that this is so, observe in general that, in the spheroid
$\dfrac{x^2}{a^2} + \dfrac{y^2 + z^2}{b^2} = 1$, the
volume included between the planes $x = \alpha$, $x = \alpha + d\alpha$, is
$= \pi(y^2 + z^2)\,d\alpha\, = \pi b^2\!\left(1 - \dfrac{\alpha^2}{a^2}\right) d\alpha$;
and thence, writing $a(1+\varepsilon)$, $b(1+\varepsilon)$ for $a$, $b$, in the shell $(a,\,b,\,b\,;\,\varepsilon)$ the volume
included between the planes $x=\alpha$, $x = \alpha + d\alpha$ is $= \pi b^2 \cdot 2\varepsilon\, d\alpha$; viz.\ this is independent
of $\alpha$, and simply proportional to $d\alpha$.  Hence, writing $b = 0$, when the shell shrinks
up into a line, the mass must be distributed uniformly over the line.  It follows that
for a line of uniform density the equipotential surfaces are each of them a prolate

\clearpage

spheroid of revolution having the extremities of the line for its foci, and that, if
we have a shell bounded by any such surface and the consecutive similar surface, with
its mass equal to that of the line, then such shell and the line will exert the same
attractions upon any point~$P$ exterior to the shell.  The attractions of the line are
most easily obtained by means of its potential~$V$; viz.\ taking $S$, $H$ for the extremities
of the line, and, as above, the origin at the middle point, and the axis of $x$ in the
direction of the line, and writing $2ae$ for the length of the line, $x$, $y$, $z$ for the
coordinates of $P$, and $r$, $s$ for the values of $HP$, $SP$
\bigl(that is, $r = \sqrt{(x - ae)^2 + y^2 + z^2}$,
$s = \sqrt{(x + ae)^2 + y^2 + z^2}$\bigr), then the potential is at once found to be
\[
V = \log\frac{x + ae + s}{x - ae + r},
\]
and we can hereby verify that the equipotential surface is in fact a spheroid of
revolution having the foci $S$, $H$; for, taking the equation of such a spheroid to be
\[
\frac{x^2}{\alpha^2} + \frac{y^2 + z^2}{\alpha^2 - a^2 e^2} = 1
\]
($\alpha$ is an arbitrary parameter, since only the value of $ae$ has been defined), we have
\[
s = \alpha + ex,\qquad r = \alpha - ex,
\]
and thence
\begin{align*}
x + ae + s &= (1 + e)(x + \alpha), \\
x - ae + r &= (1 - e)(x + \alpha),
\end{align*}
and the quotient is $= \dfrac{1 + e}{1 - e}$, a constant value, as it should be.  The equation
$V = \text{const.}$ may in fact be written
\[
\frac{1 + e}{1 - e} = \frac{x + ae + s}{x - ae + r},
\]
viz.\ this equation, apparently of the fourth order, breaks up into the twofold plane
$\dfrac{x^2}{a^2} + \dfrac{y^2 + z^2}{a^2(1-e^2)} = 1$, and the spheroid $y^2 = 0$.

The foregoing results in regard to the attraction of a line are not new.  See
Green's \textit{Essay on Electricity}, 1828, and \textit{Collected Works}, Cambridge, 1871, p.~68; also

\clearpage

Joachimsthal, ``On the Attraction of a Straight Line,'' with Sir W.\ Thomson's Note,
\textit{Camb.\ and Dubl.\ Math.\ Journ.}, vol.~\textsc{iii}. (1848), p.~93; but it does not appear to have
been noticed that they are, in fact, included in the theory of the attraction of
ellipsoids.

The like considerations show that the attractions of the ellipsoidal shell $(a,\,b,\,c\,;\,\varepsilon)$
upon an exterior point are equal to those of an elliptic disk $z = 0$, $\dfrac{x^2}{a^2-c^2} + \dfrac{y^2}{b^2-c^2} = 1$,
the mass of which is equal to that of the shell, and which has the density at the
point $(x,\,y)$ proportional to $\left(1 - \dfrac{x^2}{a^2-c^2} - \dfrac{y^2}{b^2-c^2}\right)^{-\frac{1}{2}}$.

Sir W.\ Thomson informs me that the foregoing results have long been familiar to
him.

\clearpage

% ============================================================
% PAPER 425
% Book pp.\ 34--35
% ============================================================
\section{[425]}
\subsection{Note on the Geodesic Lines on an Ellipsoid.}

\begin{center}
\textit{[From the Philosophical Magazine, vol.~\textsc{xli}. (1871), pp.~534, 535.]}
\end{center}

\textsc{The} general configuration of the geodesic lines on an ellipsoid is established by
means of the known theorem (an immediate consequence of Jacobi's fundamental formul\ae,
but which was first given by Mr Michael Roberts, \textit{Comptes Rendus}, vol.~\textsc{xxi}. p.~1470,
Dec.\ 1845) that every geodesic line touches a curve of curvature; that is, attending
to the two opposite ovals which constitute the curve of curvature, the geodesic line is in
general an infinite curve undulating between these opposite ovals, and so touching each
of them an infinite number of times (but possibly in particular cases it is a re-entrant
curve touching each oval a finite number of times).  The geodesic lines thus divide
themselves into two kinds, according as they touch a curve of curvature of the one
or the other kind; and there is besides a third limiting kind, the lines which pass
through an umbilicus: any such geodesic line passes through the opposite umbilicus,
and is in general an infinite curve passing an infinite number of times alternately
through the two umbilici; but possibly it is in particular cases a re-entrant curve
passing a finite number of times through the two umbilici.  I annex a figure giving
a general idea of the configuration of the geodesic lines drawn in different directions
from a given point~$P$ on the surface of the ellipsoid: this is drawn (as it were) on
the plane of the greatest and least axes; but it is not a perspective or geometrical
representation of any kind, but a mere diagram for the purpose in question.  We have
$A$, $A_1$, $B$, $C$, $C_1$ the extremities of the axes; $U_1$, $U_2$, $U_3$, $U_4$ the umbilici; $P$ the point
on the surface; $1P2$ and $1P4$ the curves of curvature through $P$, viz.\ these are ovals
containing the umbilici $U_1$, $U_2$ and $U_1$, $U_4$ respectively.  Then $U_1PU_3$ and $U_4PU_2$ are
the limiting geodesics passing through the umbilici; the line $TPT'$ represents a
geodesic line of the one kind, viz.\ this at~$T$ touches an oval (curve of curvature) $U_1U_4$,
and at~$T'$ the conjugate oval $U_2U_3$.  Similarly $SPS'$ is a geodesic line of the other
kind, viz.\ this at~$S$ touches an oval (curve of curvature) $U_1U_2$, and at~$S'$ the conjugate
oval $U_3U_4$; the dotted figure-of-eight curves are the loci of the points of contact
$T$, $T'$, $S$, $S'$.

\clearpage

% ============================================================
% PAPER 426
% Book p.\ 36
% ============================================================
\section{[426]}
\subsection{On a supposed New Integration of Differential Equations of the Second Order.}

\begin{center}
\textit{[From the Philosophical Magazine, vol.~\textsc{xlii}. (1871), pp.~197--199.]}
\end{center}

\textsc{This} refers to a paper, Challis, ``On the Application of a new Integration of Differential
Equations of the Second Order to some unsolved Problems in the Calculus of Variations,''
\textit{Phil.\ Mag.}\ same volume, pp.~28--40.

\clearpage

% ============================================================
% PAPER 427
% Book pp.\ 37--38
% ============================================================
\section{[427]}
\subsection{On Gauss's Pentagramma Mirificum.}

\begin{center}
\textit{[From the Philosophical Magazine, vol.~\textsc{xlii}. (1871), pp.~311, 312.]}
\end{center}

\textsc{Take} on a sphere (in the northern hemisphere) two points, $A$, $B$, whose longitudes
differ by $90^\circ$, and refer them to the equator by the meridians $AE$ and $BG$ respectively;
join $A$, $B$ by an arc of great circle, and take in the southern hemisphere the pole
$D$ of this circle; and join $D$ with $E$ and $C$ respectively by arcs of great circle.  We
have a spherical pentagon $ABCDE$, which is in fact the ``Pentagramma mirificum,''
considered by Gauss, as appearing vol.~\textsc{iii}. pp.~481--490 of the \textit{Collected Works}.  Among
its properties we have
\[
\left.
\begin{array}{l}
\text{the distance of any two non-adjacent summits} \\
\text{the inclination of any two non-adjacent sides}
\end{array}
\right\} = 90^\circ;
\]
so that each summit is the pole of the opposite side, or the pentagon is its own
reciprocal.

Each angle is the supplement of the opposite side.

If the squared tangents of the sides (or angles) taken in order are $\alpha$, $\beta$, $\gamma$, $\delta$, $\varepsilon$, then
\[
1 + \alpha = \gamma\delta, \quad 1 + \beta = \delta\varepsilon, \quad 1 + \gamma = \varepsilon\alpha,
\quad 1 + \delta = \alpha\beta, \quad 1 + \varepsilon = \beta\gamma,
\]
equivalent to three independent equations, so that any three of the quantities may be
expressed in terms of the remaining two.  (This agrees with the foregoing construction,
where the arbitrary quantities are the latitudes of $A$, $B$ respectively.)

\clearpage

Projecting from the centre of the sphere upon any plane, we have a plane
pentagon which is such that the perpendiculars let fall from the summits upon the
opposite sides respectively meet in a point.  This (as easily seen) implies that the two
portions into which each perpendicular is divided by the point in question have the
same product.

Conversely, starting from the plane pentagon, and erecting from the point of intersection
a perpendicular to the plane, the length of this perpendicular being equal to
the square root of the product in question, we have the centre of a sphere such that
the projection upon it of the plane polygon is the pentagramma mirificum.

I remark as to the analytical theory, that, taking the origin at the intersection
of the perpendiculars, and for the coordinates of the summits $(\alpha_1,\,\beta_1),\ \dots,\ (\alpha_5,\,\beta_5)$ respectively,
then we have
\[
\alpha_1\alpha_3 + \beta_1\beta_3 = \alpha_2\alpha_4 + \beta_2\beta_4 = \alpha_3\alpha_5 + \beta_3\beta_5
= \alpha_4\alpha_1 + \beta_4\beta_1 = \alpha_5\alpha_2 + \beta_5\beta_2 = -\gamma^2,
\]
where $\gamma^2$ is the above-mentioned product, or $\gamma$ is the radius of the sphere.

\begin{flushright}
\textit{Cambridge, September 14, 1871.}
\end{flushright}

\clearpage

% ============================================================
% PAPER 428
% Book p.\ 39
% ============================================================
\section{[428]}
\subsection{Note sur la Correspondance de Deux Points sur une Courbe.}

\begin{center}
\textit{[From the Comptes Rendus de l'Acad\'emie des Sciences de Paris, tom.~\textsc{lxii}.
(Janvier--Juin, 1866), pp.~586--590.]}
\end{center}

\textsc{This} is to the same effect with the paper~385, ``On the Correspondence of two Points on a Curve'';
four examples of the theory are given, the first, second, and third of them the same as in this paper---the
fourth example is as follows:

\medskip

$4^\circ$.  \textit{Recherche du nombre des points sextactiques}, c'est-\`a-dire des points qui sont tels
que par chacun passe une conique qui a dans ce point un contact du cinqui\`eme ordre
avec la courbe.  Il faut prendre pour les points $P$ les intersections avec la courbe de
la conique qui a au point $P'$ un contact du quatri\`eme ordre: les points unis seront
ceux dont il s'agit.  La courbe $\theta = 0$ est la conique qui a au point $P'$ un contact
du quatri\`eme ordre.  On a ainsi parmi les intersections le point $P'$ 5 fois; donc $k = 5$.
\`A chaque point $P'$ correspondent $2m - 5$ points~$P$; \`a chaque point~$P$, $10m^2 - 20m - 5 - 20\delta$
points~$P'$ (j'emprunte le terme $-20\delta$ d'une formule que vient de donner M.\ Zeuthen);
donc la formule donne pour le nombre des points unis
\[
10m^2 - 18m - 10 - 20\delta + 10\delta,
\]
c'est-\`a-dire
\[
10m^2 - 18m - 10 - 20\delta.
\]
Mais cette expression comprend le nombre $3m(m-2) - 6\delta$ des inflexions; en effet pour
un point d'inflexion la conique avec contact du quatri\`eme ordre se r\'eduit \`a la
tangente prise deux fois, ce qui est une conique avec contact du cinqui\`eme ordre.  Donc
enfin le nombre des points sextactiques sera
\[
m(12m - 27) - 24\delta,
\]
ou, pour une courbe sans points doubles,
\[
m(12m - 27),
\]
ce qui s'accorde avec la valeur que j'ai trouv\'ee par d'autres moyens.

\clearpage

% ============================================================
% PAPER 429
% Book p.\ 40
% ============================================================
\section{[429]}
\subsection{Sur les Coniques d\'etermin\'ees par cinq Conditions de Contact avec une Courbe donn\'ee.}

\begin{center}
\textit{[From the Comptes Rendus de l'Acad\'emie des Sciences de Paris, tom.~\textsc{lxiii}.
(Juillet--D\'ecembre, 1866), pp.~9--12.]}
\end{center}

\textsc{This} paper (dated Cambridge, 26 June 1866), contains the expressions for the numbers
$(5)$, $(4,\,1)$, $(3,\,2)$, $(3,\,1,\,1)$, $(2,\,2,\,1)$, $(2,\,1,\,1,\,1)$ and $(1,\,1,\,1,\,1,\,1)$, of the conics which satisfy
five conditions of contact with a given curve, as obtained in the paper 406 ``On the
Curves which satisfy given conditions,'' see p.~214, and which expressions were found
by the same process, viz.\ by consideration of functional equations obtained by supposing
the given curve to break up into two curves of the orders $m$ and $m'$ respectively;
there was in the expression for $(1,\,1,\,1,\,1,\,1)$ a numerical error as mentioned in the
footnote of the same page.  The paper contains also the formula
$\mu'' + \iota'' - \tau'' + \rho'' - \sigma'' = 0$, and the expression for $(2\nu,\,3\rho)$ given, pp.~203, 204.

\clearpage

% ============================================================
% PAPER 430
% Book pp.\ 41--43
% ============================================================
\section{[430]}
\subsection{Note sur quelques Formules de M.\ E.\ de Jonqui\`eres, relatives aux Courbes
qui satisfont \`a des Conditions donn\'ees.}

\begin{center}
\textit{[From the Comptes Rendus de l'Acad\'emie des Sciences de Paris, tom.~\textsc{lxiii}.
(Juillet--D\'ecembre, 1866), pp.~666--670.]}
\end{center}

\textsc{Les} formules dont il s'agit sont publi\'ees dans les \textit{Comptes Rendus}, s\'eances du 3
et du 17 septembre 1866.  En faisant une simple transformation alg\'ebrique pour y
introduire la classe $M\,(=m^2 - m)$ de la courbe donn\'ee $U^m$, et en changeant un peu la
forme, les th\'eor\`emes de M.\ de Jonqui\`eres peuvent s'\'enoncer comme il suit:

\medskip\noindent
$1^\circ$.  Le nombre des contacts des courbes $C^r$ qui ont un contact de l'ordre $n$ avec
une courbe fixe $U^m$, et qui passent en outre par $\tfrac{1}{2}r(r+3) - n$ points donn\'es, est
\[
= \tfrac{1}{2}(n+1)\bigl[nM + (2r - 2n)m\bigr].
\]

\textit{Observation}.  \'Enonc\'e de cette mani\`ere, le th\'eor\`eme s'applique m\^eme au cas $n = 0$.
En effet, pour $n = 0$, le nombre donn\'e par le th\'eor\`eme est $= mr$, qui est le nombre
des contacts de l'ordre $0$ (intersections simples) de la courbe donn\'ee $U^m$ avec une
courbe d\'etermin\'ee de l'ordre~$r$.

\medskip\noindent
$2^\circ$.  Le nombre des contacts de l'ordre $n'$ ($=$ ou $<n$) des courbes $C^r$ qui ont deux
contacts des ordres $n$ et $n'$ respectivement avec une courbe fixe $U^m$, et qui passent
en outre par $\tfrac{1}{2}r(r+3) - n - n'$ points donn\'es est
\begin{align*}
&= \tfrac{1}{4}(n+1)(n'+1)\Bigl\{[nM + (2r - 2n)m]\,[n'M + (2r - 2n')m] \\
&\qquad\qquad{} - 2(n^2 + nn' + n'^2 + n + n')\,M \\
&\qquad\qquad{} + [-4r(n + n' + 1) + 4(n^2 + nn' + n'^2 + n + n')]\,m\Bigr\}.
\end{align*}

\clearpage

\textit{Observation}.  \'Enonc\'e de cette mani\`ere, le th\'eor\`eme s'applique m\^eme aux cas
$n' = 0$, et $n' = n$.  En effet, pour $n' = 0$, le nombre donn\'e par le th\'eor\`eme est
$= (rm - n - 1) \cdot \tfrac{1}{2}(n+1)\,[nM + (2r-2n)m]$, ce qui est \'egal au nombre des courbes
qui ont avec la courbe donn\'ee $U^m$ un contact de l'ordre $n$, multipli\'e par $rm - n - 1$,
nombre des contacts de l'ordre $0$ (intersections simples) de chacune de ces courbes avec
la courbe $U^m$.  Et pour $n' = n$, le nombre des contacts est le double du nombre des
courbes~$C^r$.

Je remarque que les deux th\'eor\`emes peuvent se d\'emontrer de la mani\`ere dont
je me suis servi en cherchant le nombre des coniques qui satisfont \`a cinq conditions
donn\'ees; car, en rempla\c{c}ant la courbe $U^m$ par l'ensemble de deux courbes $m$ et $m'$,
on trouve que pour le th\'eor\`eme $1^\circ$ le nombre cherch\'e est
\[
\tfrac{1}{2}(n+1)\,[nM + (2r-2n)m] + \alpha M + \beta m,
\]
o\`u les coefficients $(\alpha,\,\beta)$ ne d\'ependent que de $(r,\,n)$; et puis, en supposant que ce
th\'eor\`eme soit connu, on trouve que pour le th\'eor\`eme $2^\circ$ le nombre cherch\'e est
\[
\tfrac{1}{4}(n+1)(n'+1)\,[nM + (2r-2n)m]\,[n'M + (2r-2n')m] + \alpha M + \beta m,
\]
o\`u de m\^eme les coefficients $(\alpha,\,\beta)$ ne d\'ependent que de $(r,\,n)$.

Or voici comment on peut d\'eterminer les coefficients dans les deux th\'eor\`emes:

Pour le th\'eor\`eme $1^\circ$, on d\'emontre que pour $U^m$ une droite, le nombre cherch\'e
est $= (n+1)(r-n)$; et que pour $U^m$ une conique, le nombre cherch\'e se d\'eduit de
l\`a en \'ecrivant $2r$ au lieu de $r$; c'est-\`a-dire, que pour la conique, le nombre est
$= (n+1)(2r-n)$.  On a donc
\begin{align*}
\alpha &+ \beta = (n+1)(r-n), \\
2\alpha &+ 2\beta = (n+1)(2r-n),
\end{align*}
et de l\`a
\[
\alpha = \tfrac{1}{2}(n+1)n, \quad \beta = \tfrac{1}{2}(n+1)(2r-2n);
\]
ce qui ach\`eve la d\'emonstration.

Pour le th\'eor\`eme $2^\circ$, on d\'emontre que pour $U^m$ une droite, le nombre cherch\'e
est $= (n+1)(n'+1)(r-n-n')(r-n-n'-1)$, et que pour $U^m$ une conique, le nombre
cherch\'e se d\'eduit de l\`a en \'ecrivant $2r$ au lieu de $r$; c'est-\`a-dire, pour la conique,
le nombre est
\[
= (n+1)(n'+1)(2r-n-n')(2r-n-n'-1).
\]
On a donc
\begin{align*}
(n+1)(n'+1)(r-n-n')(r-n-n'-1) &= (n+1)(n'+1)(r-n)(r-n') + \alpha + \beta, \\
(n+1)(n'+1)(2r-n-n')(2r-n-n'-1) &= (n+1)(n'+1)(2r-n)(2r-n') + \alpha + 2\beta,
\end{align*}
cela donne pour $\alpha$ et $\beta$ les valeurs
\begin{align*}
\alpha &= \tfrac{1}{4}(n+1)(n'+1)\,[-2(n^2 + nn' + n'^2 + n + n')], \\
\beta  &= \tfrac{1}{4}(n+1)(n'+1)\,[-4r(n + n' + 1) + 4(n^2 + nn' + n'^2 + n + n')],
\end{align*}
et la d\'emonstration est ainsi achev\'ee.

\clearpage

Je remarque que sous les formes ici donn\'ees les deux th\'eor\`emes s'appliquent \`a
une courbe $U^m$ avec des points doubles, mais sans point de rebroussement.

Le th\'eor\`eme dont je me suis servi pour la d\'etermination des coefficients peut
s'\'enoncer sous la forme plus g\'en\'erale que voici, savoir:

En d\'enotant par $\phi(r,\,n,\,n',\dots)$ le nombre des courbes $C^r$ qui ont avec une droite
donn\'ee des contacts des ordres $n$, $n'$, \dots\ et qui passent en outre par $\tfrac{1}{2}r(r+3) - n - n'\dots$
points donn\'es, alors si, au lieu de la droite donn\'ee, on a une conique donn\'ee, le
nombre des courbes $C^r$ sera $= \phi(2r,\,n,\,n',\dots)$.

En effet, l'\'equation de la courbe cherch\'ee $C^r$ contient des coefficients ind\'etermin\'es,
lesquels, par les conditions de passer par les points donn\'es, se r\'eduisent lin\'eairement
\`a $n + n' + \ldots + 1$ coefficients; en d\'enotant par $(A,\,B,\,\dots)$ ces coefficients, l'\'equation de la
courbe contiendra lin\'eairement $(A,\,B,\dots)$ et sera ainsi de la forme $(A,\,B,\dots\,\,\check{a}\,\,x,\,y,\,z)^r = 0$.
L'\'equation de la droite donn\'ee est satisfaite en prenant pour $(x,\,y,\,z)$ des fonctions
lin\'eaires d\'etermin\'ees d'un param\`etre variable~$\theta$; donc, en coupant la courbe $C^r$ par
la droite donn\'ee, on obtient une \'equation $(A,\,B,\dots\,\,\check{a}\,\,\theta,\,1)^r = 0$, et en exprimant que
cette \'equation ait $n$ racines \'egales, $n'$ racines \'egales, etc., on obtient entre $(A,\,B,\,C,\dots)$
des \'equations, lesquelles, en \'eliminant tous les coefficients, except\'es deux quelconques
$(A,\,B)$, conduisent \`a une \'equation finale $(A,\,B)^p = 0$, et le degr\'e $p$ de cette \'equation
est ce qu'il s'agissait de trouver, le nombre des courbes $C^r$.  Si au lieu d'une droite
donn\'ee on a une conique donn\'ee, il n'y a rien \`a changer, sinon que les coordonn\'ees
$(x,\,y,\,z)$ doivent \^etre remplac\'ees par des fonctions quadratiques de~$\theta$; on a ainsi une
\'equation $(A,\,B,\dots\,\,\check{a}\,\,\theta,\,1)^{2r} = 0$, qui conduit \`a une \'equation finale $(A,\,B)^{p'} = 0$, o\`u $p'$
est la m\^eme fonction de $(2r,\,n,\,n',\dots)$ qu'est $p$ de $(r,\,n,\,n',\dots)$; et le nombre des
courbes $C^r$ est $= p'$.  Le th\'eor\`eme est donc d\'emontr\'e.  Et, pr\'ecis\'ement de la m\^eme
mani\`ere, on d\'emontre le th\'eor\`eme encore plus g\'en\'eral:

En d\'enotant par $\phi(r,\,n,\,n',\dots)$ le nombre des courbes $C^r$ qui ont avec une droite
donn\'ee des contacts des ordres $n$, $n'$, \dots, et qui passent en outre par $\tfrac{1}{2}r(r+3) - n - n'\dots$
points donn\'es, alors si, au lieu de la droite donn\'ee, on a une courbe unicursale donn\'ee
de l'ordre $m$, le nombre des courbes $C^r$ est $= \phi(mr,\,n,\,n',\dots)$.

On aurait pu se servir directement de cela pour d\'emontrer les th\'eor\`emes $1^\circ$ et $2^\circ$.
Par exemple, pour le th\'eor\`eme $1^\circ$, la consid\'eration de la courbe unicursale $U^m$ donne
\[
\alpha M + \beta m = \alpha(2m - 2) + \beta m = (n+1)(mr - n),
\]
c'est-\`a-dire
\[
\alpha = \tfrac{1}{2}(n+1)n, \quad \beta = \tfrac{1}{2}(n+1)(2r - 2n),
\]
comme auparavant.

\clearpage

% ============================================================
% PAPER 431
% Book pp.\ 44--46
% ============================================================
\section{[431]}
\subsection{Sur la Transformation Cubique d'une Fonction Elliptique.}

\begin{center}
\textit{[From the Comptes Rendus de l'Acad\'emie des Sciences de Paris, tom.~\textsc{lxiv}.
(Janvier--Juin, 1867), pp.~560--563.]}
\end{center}

\textsc{Soit} $U = (a,\,b,\,c,\,d,\,e\,\,\check{a}\,\,x,\,1)^4$ une fonction quartique quelconque de~$x$: $I$, $J$ les
deux invariants:
\[
(I = ae - 4bd + 3c^2,\quad J = ace - ad^2 - b^2e + 2bcd - c^3),
\]
et prenons $\Omega = \dfrac{I^3 - 27J^2}{I^3}$ pour l'invariant absolu de~$U$.  Soient de m\^eme $U' = (a',\,b',\,\dots\,\,\check{a}\,\,x,\,1)^4$
et $\Omega' = \dfrac{I'^3 - 27J'^2}{I'^3}$ l'invariant absolu de~$U'$.  En supposant que $\sqrt[4]{U}$, $\sqrt[4]{U'}$ soient les
radicaux des deux fonctions elliptiques li\'ees par la transformation du troisi\`eme ordre
ou cubique, on peut se proposer la question quelle est la relation entre les deux
invariants absolus $\Omega$, $\Omega'$?  J'ai trouv\'e cette relation d'abord par des consid\'erations
g\'eom\'etriques qui me furent sugg\'er\'ees par une lettre de M.\ Sylvestre; puis je l'ai d\'eduite
des formules pour la transformation cubique donn\'ees par M.\ Hermite, \textit{Crelle}, t.~\textsc{lx}., 1862,
p.~304), et enfin, \`a l'aide d'une consid\'eration tir\'ee de ces formules, j'ai r\'eussi \`a
l'obtenir au moyen des formules des \textit{Fundamenta Nova}.  Je vais donner ici cette derni\`ere
investigation de la relation dont il s'agit.

En supposant que les fonctions $U$, $U'$ soient transform\'ees lin\'eairement en
$(1 - x^2)(1 - k^2x^2)$, $(1 - y^2)(1 - \lambda^2y^2)$ respectivement, pour exprimer la liaison entre les
modules $k$, $\lambda$, au lieu de l'\'equation explicite entre $\sqrt[4]{k}$, $\sqrt[4]{\lambda}$ (\textit{Fund.\ Nova}, p.~23), je
me sers des formules, p.~25, lesquelles en y \'ecrivant
\[
-\beta = \frac{\alpha + 2}{2\alpha + 1},
\]

\clearpage

c'est-\`a-dire
\[
2\alpha\beta + \alpha + \beta = -2,
\]
deviennent
\[
k^2 = -\alpha^2\beta,\quad \lambda^2 = \alpha\beta^2.
\]

Les transformations lin\'eaires donnent sans peine
\[
\Omega = \frac{108k^2(k^2 - 1)^2}{(k^4 + 14k^2 + 1)^3},
\quad
\Omega' = \frac{108\lambda^2(\lambda^2 - 1)^2}{(\lambda^4 + 14\lambda^2 + 1)^3},
\]
et il s'agit entre ces \'equations d'\'eliminer $\alpha$, $\beta$, $k$, $\lambda$ de mani\`ere \`a obtenir une \'equation
entre $\Omega$, $\Omega'$.

J'\'ecris
\[
\Omega = \frac{\tfrac{1}{4}(2\alpha + 1)(\alpha + 2)(\alpha - 1)^4}{(\alpha^2 + 4\alpha + 1)^3},
\quad
\Omega' = \frac{\tfrac{1}{4}(2\beta + 1)(\beta + 2)(\beta - 1)^4}{(\beta^2 + 4\beta + 1)^3}.
\]

L'\'equation entre $\alpha$, $\beta$ donne
\[
2\beta + 1 = \frac{-3(\alpha + 1)}{2\alpha + 1},\quad
\beta + 2 = \frac{3(\alpha^2 + 4\alpha + 1)}{2\alpha + 1},\quad
\beta - 1 = \frac{-3(\alpha + 1)^2}{2\alpha + 1},
\quad
\beta^2 + 4\beta + 1 = \frac{3(\alpha^2 + 4\alpha + 1)}{(2\alpha + 1)^2};
\]
et on a de l\`a
\[
\beta' = \frac{\beta^2(\beta - 1)^2(\beta + 1)^2}{\beta^2(\beta + 2)^2}
\]
puis, en faisant attention \`a l'identit\'e
\[
(2\alpha + 1)(\alpha + 2)(\alpha - 1)^4 + 27\alpha(\alpha + 1)^4 = 2(\alpha^2 + 4\alpha + 1)^3,
\]
on obtient entre $\alpha'$, $\beta'$, la relation tr\`es simple $\alpha' + \beta' = 1$.

L'expression de $k^2$ donne
\begin{align*}
k^2 &= \frac{\alpha^2(\alpha + 2)}{2\alpha + 1}, \\
k^2 - 1 &= \frac{(\alpha - 1)(\alpha + 1)^2}{2\alpha + 1}, \\
k^4 + 14k^2 + 1 &= \frac{1}{(2\alpha + 1)^2}\bigl\{\alpha^4(\alpha + 2)^2 + 14\alpha^2(\alpha + 2)(2\alpha + 1) + (2\alpha + 1)^2\bigr\}, \\
&= \frac{1}{(2\alpha + 1)^2}\bigl\{(\alpha^2 + 4\alpha + 1)(\alpha^4 + 3\alpha^3 + 16\alpha^2 + 3\alpha + 1)\bigr\},
\end{align*}
et on a de l\`a
\[
\Omega = \frac{108\alpha^4(2\alpha + 1)(\alpha + 2)(\alpha - 1)^2(\alpha + 1)^4}{(\alpha^2 + 4\alpha + 1)^3 \cdot (\alpha^4 + 3\alpha^3 + 16\alpha^2 + 3\alpha + 1)^3}.
\]

\clearpage

Or on a
\begin{align*}
\alpha' - 1 &= \frac{\tfrac{1}{4}(2\alpha + 1)(\alpha + 2)(\alpha - 1)^4 - (\alpha^2 + 4\alpha + 1)^3}{(\alpha^2 + 4\alpha + 1)^3}
= \frac{-\tfrac{27}{4}\alpha(\alpha + 1)^4}{(\alpha^2 + 4\alpha + 1)^3}, \\
8\alpha' + 1 &= \frac{4(2\alpha + 1)(\alpha + 2)(\alpha - 1)^4 + (\alpha^2 + 4\alpha + 1)^3}{(\alpha^2 + 4\alpha + 1)^3}
= \frac{9(\alpha^4 + 3\alpha^3 + 16\alpha^2 + 3\alpha + 1)^2}{(\alpha^2 + 4\alpha + 1)^3},
\end{align*}
et de l\`a, en formant l'expression de la fonction $\dfrac{-64\alpha'(\alpha' - 1)^2}{(8\alpha' + 1)^3}$, on la trouve \'egale \`a
la valeur qui vient d'\^etre donn\'ee pour $\Omega$ en termes de $\alpha$: on a donc
\[
\Omega = \frac{-64\alpha'(\alpha' - 1)^2}{(8\alpha' + 1)^3}
\]
et de m\^eme
\[
\Omega' = \frac{-64\beta'(\beta' - 1)^2}{(8\beta' + 1)^3}.
\]

Avec la relation $\alpha' + \beta' = 1$, l'\'elimination de $\alpha'$, $\beta'$ entre ces \'equations ne pr\'esente
pas de difficult\'e.

\clearpage

% ============================================================
% PAPER 432
% Book p.\ 47
% ============================================================
\section{[432]}
\subsection{Th\'eor\`eme relatif \`a la Th\'eorie des Substitutions.}
\begin{center}
\textit{Extrait d'une lettre adress\'ee \`a M.\ J.\ A.\ Serret.}
\end{center}

\begin{center}
\textit{[From the Comptes Rendus de l'Acad\'emie des Sciences de Paris, tom.~\textsc{lxvii}.
(Juillet--D\'ecembre, 1868), pp.~784, 785.]}
\end{center}

\textsc{On} peut \'enoncer par rapport aux substitutions un th\'eor\`eme qui comprend les
trois th\'eor\`emes \textsc{iii}., \textsc{iv}., \textsc{v}., t.~\textsc{ii}.\ pp.~260--263 de votre \textit{Cours d'Alg\`ebre Sup\'erieure}.

Pour un nombre quelconque $\mu$ on peut former avec les $\phi(\mu)$ nombres inf\'erieurs
et premiers \`a~$\mu$ plusieurs syst\`emes de nombres lesquels sont chacun un syst\`eme
conjugu\'e $(\mathrm{mod.}\ \mu)$; c'est-\`a-dire que le produit de deux nombres quelconques d'un tel
syst\`eme est congru suivant le module~$\mu$, \`a un nombre du syst\`eme.  Comme cas
extr\^emes, l'unit\'e est un tel syst\`eme, et les $\phi(\mu)$ nombres forment aussi un syst\`eme
conjugu\'e.

Pour $\mu$ premier, en d\'enotant par $a$ une racine primitive de~$\mu$ et par~$f$ un
diviseur quelconque de $\mu - 1$, les nombres $\equiv a^{f x}\ (\mathrm{mod.}\ \mu)$, $x$ \'etant un entier quelconque,
forment un syst\`eme conjugu\'e.  Et g\'en\'eralement pour un nombre $\mu$ quelconque ce
nombre a des racines quasi-primitives $\alpha$, $\beta$, $\gamma$, \dots, aux exposants $A$, $B$, $C$, \dots, tels que
$\alpha^A \equiv 1\ (\mathrm{mod.}\ \mu)$, $\beta^B \equiv 1\ (\mathrm{mod.}\ \mu)$, \dots\ et $ABC \dots = \phi(\mu)$.  En choisissant une combinaison
quelconque, par exemple $\alpha$, $\beta$ de ces racines, soient $f$, $g$ des diviseurs quelconques de $A$, $B$
respectivement, les nombres $\equiv \alpha^{f x}\beta^{g y}\ (\mathrm{mod.}\ \mu)$ forment un syst\`eme conjugu\'e, l'ordre du
syst\`eme ou nombre des termes \'etant $= \dfrac{AB}{fg}$.

Cela \'etant, on a ce th\'eor\`eme:

\medskip
\textit{Une substitution~$T$ quelconque de l'ordre~$\mu$ \'etant form\'ee avec $n$ lettres, l'on forme
toutes les substitutions~$S$ telles que le produit $STS^{-1}$ se r\'eduise \`a une puissance de~$T$
dont l'exposant soit un nombre quelconque appartenant \`a un syst\`eme conjugu\'e $(\mathrm{mod.}\ \mu)$,
les substitutions~$S$ constitueront un syst\`eme conjugu\'e de l'ordre $\theta M$, o\`u $\theta$ d\'enote l'ordre
du syst\`eme conjugu\'e $(\mathrm{mod.}\ \mu)$ et $M$ le nombre des substitutions \'echangeables avec~$T$.}

\medskip
La d\'emonstration est tout \`a fait la m\^eme que celle que vous donnez p.~62 pour
votre th\'eor\`eme \textsc{iv}, en y ajoutant seulement que les nombres $i$, $j$ qui appartiennent
au syst\`eme conjugu\'e $(\mathrm{mod.}\ \mu)$ auront leur produit $ij$ congru \`a un nombre de ce m\^eme
syst\`eme conjugu\'e.

\clearpage

% ============================================================
% PAPER 433
% Book pp.\ 48--50
% ============================================================
\section{[433]}
\subsection{Sur les Surfaces T\'etra\'edrales.}

\begin{center}
\textit{[Notes to the work De la Gournerie, ``Recherches sur les surfaces r\'egl\'ees t\'etra\'edrales
sym\'etriques,'' 8vo.\ Paris, 1867.]}
\end{center}

\begin{center}
\textsc{Premier M\'emoire.\ \ Notes pp.~190--193.}
\end{center}

\medskip\noindent
\textit{\'Equations de certaines}
\medskip

$1^\circ$.  L'\'equation
\begin{align*}
0 = {}& a^4 f^4 v^4 + b^4 g^4 u^4 + c^4 h^4 w^4 \\
{}+{}& 2c^2 ag\,(af - bg)\,y^2 z^2 - 2c^2 bf\,(af - bg)\,x^2 z^2 \\
{}-{}& 2b^2 ah\,(ch - af)\,z^2 x^2 + 2b^2 ch\,(ch - af)\,y^2 z^2 \\
{}+{}& 2a^2 gh\,(bg - ch)\,w^2 x^2 + 2a^2 bh\,(bg - ch)\,z^2 y^2 \\
{}+{}& 2g^2 hf\,(ch - af)\,w^2 y^2 \\
{}+{}& a^2\,(b^2 y^2 + c^2 h^2 - 4bg\,ch)\,y^2 z^2
       + b^2\,(c^2 h^2 + a^2 f^2 - 4ch\,af)\,z^2 x^2 \\
{}+{}& g^2\,(c^2 h^2 + a^2 f^2 - 4ch\,af)\,w^2 y^2 \\
{}+{}& f^2\,(b^2 y^2 + c^2 h^2 - 4bg\,ch)\,w^2 x^2 \\
{}+{}& 2bf\,(af bg + c^2 h^2 - 2ch\chi)\,x^2 z^2 w^2 \\
{}-{}& 2ag\,(af bg + c^2 h^2 - 2ch\chi)\,y^2 w^2 z^2 \\
{}+{}& 2ah\,(ch af + b^2 y^2 - 2bg\chi)\,z^2 h u^2
       - 2bh\,(bg ch + a^2 f^2 - 2af\chi)\,z^2 \chi w^2 \\
{}-{}& 2gh\,(bc gh + af^2 - 2af\chi)\,w^2 y^2 z^2
       - 2hf\,(c a hf + b y^2 - 2bg\chi)\,w^2 z^2 u^2 \\
{}+{}& 2\Omega h f^2 z^2 w^2.
\end{align*}

\medskip\noindent
\textsc{Surfaces du Huiti\`eme Ordre.}
\medskip

\begin{align*}
{}+{}& 2b^2 cf\,(ch - af)\,x^2 z^2 - 2bc f^2\,(bg - ch)\,x^2 w^2 \\
{}-{}& 2a^2 cg\,(bg - ch)\,y^2 z^2 - 2c a g^2\,(ch - af)\,y^2 w^2 \\
{}-{}& 2abh^2\,(af - bg)\,z^2 w^2 + 2h^2 f g\,(af - bg)\,w^2 z^2 \\
{}-{}& c^2\,(ch - af)\,z^2 x^2 + c^2\,(a^2 f^2 + b^2 g^2 - 4af bg)\,x^4 y^4 \\
{}-{}& 4\,(ch - af)\,w^2 y^2 + h^2\,(a^2 f^2 + b^2 g^2 - 4af bg)\,w^2 z^2 \\
{}-{}& 2cf\,(ch af + b^2 f^2 - 2bg\chi)\,x^2 w^2 y^2 + 2bc\,(bc gh + a^2 f^2 - 2af\chi)\,x^2 y^2 z^2 \\
{}+{}& 2cg\,(bg ch + a^2 f^2 - 2af\chi)\,y^2 x^2 w^2 + 2ca\,(ca hf + by^2 - 2bg\chi)\,y^2 z^2 x^2 \\
{}+{}& 2ab\,(ab fg + c^2 h^2 - 2ch\chi)\,z^2 x^2 y^2 \\
{}-{}& 2\,(ab fg + c^2 h^2 - 2ch\chi)\,w^2 x^2 y^2.
\end{align*}

\clearpage

(o\`u, pour abr\'eger, on a \'ecrit $\chi = af + bg + ch$, et $\Omega$ est une quantit\'e quelconque) est
celle d'une surface du huiti\`eme ordre qui a pour courbes doubles les quatre coniques
\begin{align*}
x &= 0, & {} +c y^2 - b z^2 + f w^2 &= 0; \\
y &= 0, & -c x^2 \phantom{+ay^2} + a z^2 + g w^2 &= 0; \\
z &= 0, & {} +b x^2 - a y^2 \phantom{+gw^2} + h w^2 &= 0; \\
w &= 0, & -f x^2 - g y^2 - h z^2 &= 0.
\end{align*}

En d\'eterminant $\Omega$, savoir, en \'ecrivant $\lambda + \mu + \nu = 0$, $af\lambda^2 + bg\mu^2 + ch\nu^2 = 0$ (ce qui
donne deux syst\`emes de valeurs de $\lambda : \mu : \nu$), et puis
\[
\Omega = 4 \sum \frac{\mu^2 \nu^2}{\lambda^2}\,(af + bg)^2\,ch
- 2 \sum \frac{\mu^2 \nu^2}{\lambda^2}\,(af + bg)^2
- 4\,(af - bg)\,(bg - ch)\,(ch - af),
\]
la surface devient une surface r\'egl\'ee, savoir, la quadrispinale de M.\ de la Gournerie;
et, en particulier, en supposant $\dfrac{1}{af} + \dfrac{1}{bg} + \dfrac{1}{ch} = 0$, on obtient
\[
\Omega = (-2 - 4) = -6\,(af - bg)\,(bg - ch)\,(ch - af),
\]
et la surface sera d\'eveloppable.

\medskip\noindent
$2^\circ$.  L'\'equation
\begin{align*}
0 = {}& a^2 y^2 z^4 + b^2 z^2 x^4 + c^2 x^2 y^4 + f^2 z^4 w^4 + g^2 x^4 w^4 + h^2 y^4 w^4 \\
{}+{}& 2bf x^2 z^2 w^2 - 2c f z x^2 y^2 w^2 + 2bc z x^2 y^2 z^2 \\
{}-{}& 2ag y^2 z^2 w^2 + \ldots + 2c g x^2 y^2 w^2 + 2 c a x^2 y^2 z^2 \\
{}+{}& 2ah z^2 y^2 w^2 - 2bh z^2 x^2 w^2 + \ldots + 2ab x^2 y^2 z^2 \\
{}-{}& 2gh w^2 y^2 z^2 - 2hf w^2 z^2 x^2 - 2fg w^2 x^2 y^2 \\
{}+{}& 2\Omega\, x^2 y^2 z^2 w^2
\end{align*}
(o\`u $\Omega$ est une quantit\'e quelconque) est celle d'une surface du huiti\`eme ordre ayant
pour courbes doubles les quatre courbes du quatri\`eme ordre
\begin{align*}
x &= 0, & -h z^2 w^2 + g w^2 y^2 - a y^2 z^2 &= 0; \\
y &= 0, & -h z^2 w^2 + f w^2 x^2 + b z^2 x^2 &= 0; \\
z &= 0, & {} + g w^2 x^2 - f w^2 y^2 + c x^2 y^2 &= 0; \\
w &= 0, & {} - a y^2 z^2 - b z^2 x^2 - c x^2 y^2 &= 0.
\end{align*}

En \'ecrivant
\[
\lambda + \mu + \nu = 0, \quad af\lambda^2 + bg\mu^2 + ch\nu^2 = 0,
\]
\dots

\end{document}
